%% response_TCAD-2026-0901.tex — point-by-point response to the reviewers
%% Compile: pdflatex response_TCAD-2026-0901
%% (supersedes response.tex, which was drafted before the revision landed)

\documentclass[11pt]{article}
\usepackage[a4paper,margin=1in]{geometry}
\usepackage{xcolor}
\usepackage{enumitem}
\usepackage{booktabs}
\usepackage{amsmath}
\usepackage{hyperref}
\hypersetup{colorlinks=true,linkcolor=blue,urlcolor=blue}

\newcommand{\dSABRE}{d\textsc{sabre}}
\newcommand{\TeleSABRE}{Tele\textsc{sabre}}
\newcommand{\SABRE}{\textsc{sabre}}

\newenvironment{comment}{\medskip\noindent\itshape}{\par}
\newcommand{\response}{\smallskip\noindent\textbf{Response.}\ }
\newcommand{\where}[1]{\textcolor{blue}{\textbf{Changes:} #1}}

\setlength{\parskip}{2pt}
\emergencystretch=2em

\title{Response to Reviewers\\[3pt]
\large Manuscript TCAD-2026-0901\\
``\dSABRE{}: A \SABRE{}-Style Router for Multi-Core Distributed Quantum
Computers''}
\author{Sanjiang Li}
\date{}

\begin{document}
\maketitle

\noindent
Dear Editors and Reviewers,

\medskip\noindent
Thank you for the careful and constructive reviews. The shared thrust of the
comments---that the design read as a collection of heuristics, that key
choices and cross-model comparisons needed justification, and that the
exposition needed reorganisation---led me to revise the paper substantially
rather than incrementally. Because the rewrite touches almost every section,
colour-marking the changes would have marked most of the manuscript; instead
each response below points to the sections, equations, and tables that carry
it. Section, equation, and table numbers refer to the revised manuscript.

\medskip\noindent
Two of the reviewers' points are answered by argument and scope rather than by
new experiments, and I say so plainly where that is the case: the novelty
question (Comment~2.1) and the request for latency and fidelity metrics
(Comment~2.4(d)). Everything else is answered with new text, new theory, or
new measurements.

\section*{Summary of major changes}

\begin{enumerate}[leftmargin=*,itemsep=3pt]

\item \textbf{A unifying principle, stated and enforced.} Every inter-core
move is now formulated and scored as one \emph{macro-action}
(Eq.~3): the resource it spends, the capacity it consumes, and the progress it
buys. The two cost terms beyond \SABRE{}'s are derived, not added: port
availability and destination capacity are exactly the two routing-state
constraints a \emph{directional} teleport introduces that a symmetric,
occupancy-preserving SWAP does not (Sections~I and~III-A).

\item \textbf{Mechanisms removed, not accumulated.} The hop-gain reward, the
proactive congestion-relief pass, and the \SABRE{}-inherited node-decay
penalty of the original submission did not follow from that principle and were
not reliably load-bearing when re-measured; all three are removed. What
remains is the score, a legality condition, and a recovery transaction, and
Table~IV quantifies what each contributes.

\item \textbf{Completion is proven rather than hoped for.} The abort-prone
checkpoint escape of the original (up to 50 failed recovery attempts, after
which the router gave up) is replaced by a legality floor (Eq.~4), checkable
pre-routing capacity conditions (Eqs.~12--14), and a guaranteed recovery
transaction (Algorithm~2), with Lemmas~1--2 and Theorem~1 (termination with
success). The full proof is in the supplementary material.

\item \textbf{Evaluation broadened.} The 64-qubit suite grows from six to nine
circuit families (adding QPEexact, QAOA, and a $13\,040$-CX 16-bit
Multiplier); two IBM heavy-hex core networks (a ring and a star) are added
alongside the two grid-of-grids devices; and the scalability study covers two
families (QFT, QPEexact) to 360 logical qubits---45 circuit--architecture
instances in total.

\item \textbf{Cross-model comparisons made controlled.} \TeleSABRE{} is fixed
as the like-for-like baseline; the \texttt{pytket-dqc} comparison is reframed
as a cost-model study with explicit capacity and entanglement-lifetime audits;
and a DMapS comparison with its device-model assumptions spelled out is added
(Sections~IV-E and~V).

\item \textbf{Exposition reorganised.} Rewritten abstract and introduction,
notation consolidated in Table~I with defaults, rewritten figure captions,
consistent naming, and---for reproducibility---the machine, implementation
language, and tool versions now stated in Section~IV-A.

\end{enumerate}

\section*{Note on the reported numbers, which differ from the original
submission}

The reviewers will reasonably diff the two versions, so I account for the
differences up front. There are four independent causes---(a), (b), (d) and (e) below---one of them
a correction to my own experimental driver that made the \emph{baseline}
stronger; paragraph~(c) unpicks the one headline number in which the first two
combine.

\paragraph{(a) The algorithm changed.}
Items 1--3 above, together with an initial-layout stage that now reserves
capacity in every core by construction (Section~III-G), move the headline
geometric-mean EPR reductions against \TeleSABRE{} from $41.1/44.1/43.9\%$ to
$50.8/62.0/49.1\%$ on the 25-, 36-, and 64-qubit suites.

\paragraph{(b) A misspelled baseline parameter has been corrected, which
mostly helped the baseline.}
The original submission's \TeleSABRE{} runs passed the key
\texttt{init\_layout\_hun\_}\allowbreak\texttt{free\_qubit}, a misspelling of
\TeleSABRE{}'s \texttt{init\_layout\_hun\_}\allowbreak\texttt{min\_free\_qubit};
the misspelled key was silently
ignored and the parameter fell back to the tool's internal default rather than
the value in \TeleSABRE{}'s own shipped configuration. The revision passes the
correct key. This changes the baseline's counts, and on the larger circuits it
improves them:

\begin{center}
\begin{tabular}{@{}lrrl@{}}
\toprule
64q circuit & submitted & revised & effect on the baseline \\
\midrule
GHZ        &  11 &  16 & worse \\
Graphstate &  76 &  62 & better \\
QFT        & 410 & 312 & better \\
\bottomrule
\end{tabular}
\end{center}

\noindent
The revised figures are reproduced exactly by re-running \TeleSABRE{}
unchanged with the corrected key, so the difference is attributable to that
key alone and not to any other change on our side. Net, the corrected baseline
is the stronger one on the circuits that dominate the geometric mean, so the
reported margins are more conservative than the original's, not less.

\paragraph{(c) The 64-qubit GHZ regression is gone, for two reasons that
should be separated.}
The original reported $+45.5\%$ on 64-qubit GHZ; the revision reports
$-62.5\%$. Part of that is (b): the baseline's count moves from 11 to 16.
The rest is a genuine improvement in \dSABRE{}, whose own GHZ count falls from
16 to 6 EPR pairs. The cause is the initial-layout fix, not the score: the
submitted layout removed the four globally lowest-degree sites of the whole
architecture graph, which under core-major numbering are always the four
corners of core~0, so cores $1,\dots,K-1$ received \emph{no} reserved escape
capacity and regularly came out fully packed---on the 64-qubit H-grid, GHZ
landed with four of six cores at $100\%$ occupancy. Section~III-G now reserves
the most remote low-degree sites of \emph{every} core, and the legality floor
of Eq.~(4) keeps ordinary routing from filling a core thereafter. A star-shaped
interaction pattern such as GHZ is precisely the case that punishes the old
rule, which is why Reviewer~1 was right to read the regression as a signal
rather than noise.

\paragraph{(d) Compile time improved, and the accounting became stricter.}
The original submission reported \dSABRE{} as $5$--$12\times$ slower than
\TeleSABRE{} and attributed this to the pure-Python versus compiled-C\texttt{++}
constant factor \emph{and} to per-iteration NetworkX shortest-path lookups that
\TeleSABRE{} caches. The second half of that diagnosis was the actionable one
and has been addressed: physical distances are no longer answered by
per-iteration graph searches but composed on demand through the port
decomposition of Eq.~(1) and memoised, and the router's DAG bookkeeping
(remaining in-degrees, gate counters, checkpoint/rollback, and the incremental
front-layer and lookahead score deltas) is now maintained incrementally rather
than recomputed. This rewrite was verified to be output-identical to the
previous implementation over 207 routing passes across the 25-, 36-, 64-, and
360-qubit suites, at $3.2$--$5.6\times$ less compile time; the remaining
speed-up comes from the removed mechanisms. Measured end-to-end, the
geometric-mean compile time on the 64-qubit suite is now $1.9$\,s for
\dSABRE{} against $4.7$\,s for \TeleSABRE{}, and $23.7$\,s against $136.3$\,s
on the largest circuit. Two points of honesty about that comparison:
\dSABRE{} remains pure Python and \TeleSABRE{} a compiled C\texttt{++} binary
(now stated in Section~IV-A, with the machine), and \dSABRE{}'s reported time
now covers \emph{more} work than the original's---layout generation and all
nine routing passes (three \texttt{SabreLayout} seeds $\times$ three passes),
where the original timed three passes from one seed.

\paragraph{(e) The 360-qubit device was corrected, and is now harder.}
The original reported the 360-qubit QFT in $801$\,s. While regenerating the
series I found that the scalability driver had been wired to a six-core device
of 81 qubits per core (core-graph diameter 3, 486 physical qubits), which is
not the device the scalability text describes. The revision uses the twenty
$5\times5$-core device the text specifies (diameter 7, 500 physical qubits) and
completes the same circuit in $48$\,s. The two timings are therefore not on
identical hardware, and the corrected device is the more demanding of the two.

\medskip\noindent
Per-seed counts, benchmark provenance, the parameter and cost-model sweeps, the
completion proof, and the capacity, lifetime, and action-level audits behind
the cross-model comparisons are in the supplementary material, which is
included with this submission and also released with the code.

\clearpage
\section*{Reviewer 1}

\begin{comment}
Comment 1.1. The overall design appears to be primarily driven by a collection
of heuristic optimizations addressing several observed failure modes rather
than being derived from a unified optimization principle. The interaction
between these mechanisms is not clearly justified in terms of a consistent
global objective; the method reads more like a set of incremental fixes.
\end{comment}

\response
I agree, and this comment shaped the revision. The design is now derived from
one principle: an inter-core move is a single macro-action, scored by
Eq.~(3)---resource spent $+$ capacity consumed $+$ progress bought---where the
two non-\SABRE{} terms are precisely the two routing-state constraints a
directional teleport introduces that a symmetric, occupancy-preserving SWAP
does not (Sections~I, III-A). Three mechanisms that did not follow from this
principle were removed rather than defended: the hop-gain reward, proactive
congestion relief, and the node-decay penalty. Feasibility is now deliberately
separated from ranking: the legality floor (Eq.~4) is a hard admission test
applied at candidate generation---``a weight can be outvoted by a large
progress term, an admission test cannot''---and recovery is a proven
transaction (Algorithm~2, Theorem~1) rather than a third heuristic action. The
redesigned ablation supports exactly this decomposition: the graded capacity
score buys efficiency ($+9.4\%$ EPR when removed); the floor is what the
completion proof requires, and since the score's penalty band contains the
region the floor forbids, the two overlap in practice---removing either alone
is survivable, removing both costs $+87.8\%$ \emph{and} aborts two circuits;
and the BFS depth-ordered extended set contributes $+12.4\%$ (Table~IV). The Conclusion
states the claim directly: the advance is not the number of heuristic terms
but the separation of resource-aware ranking from an explicit legality
condition and a completion-preserving recovery.

\where{Abstract; Section~I (design-principle and scope paragraphs);
Section~III-A (Eq.~3 and the legality rule); Section~III-F; Table~IV;
Section~VII.}

\begin{comment}
Comment 1.2. The exclusive use of teledata, while \TeleSABRE{} supports both
teledata and telegate, is not sufficiently motivated. This raises questions
about fairness and completeness of the comparison.
\end{comment}

\response
The scope is now both motivated and measured. Teledata is chosen because each
EPR pair is consumed immediately, which makes the EPR accounting independent
of entanglement-storage lifetime and isolates the capacity-aware routing
problem (Section~I, ``Scope of communication''). Telegate is a distinct joint
primitive-selection problem, not an omitted zero-cost candidate---and rather
than assert that, Section~VI now quantifies it. Implementing a telegate
candidate inside the same score (it consumes no destination capacity, so
$c_{\mathrm{cap}}=0$; it needs staging for two link qubits; its progress term
is the front-layer gate it retires outright) and letting it compete with
teledata costs $93.5\%$ \emph{more} EPR on the 64-qubit suite, and the best of
seventeen scoring variants swept is still $1.0\%$ worse than teledata alone.
The reason is amortisation, and it is measurable: one teledata pair already
serves $15.75$ gates on this suite on average (median $8$), with $83.4\%$ of
relocations serving two or more, whereas an ungrouped telegate serves exactly
one. A telegate pays
only in grouped form, which requires commutation-derived burst extraction,
choice of a beneficiary core, and entanglement held for the group's
duration---identified as future work with the relevant literature.

On fairness: the \TeleSABRE{} comparison charges \TeleSABRE{} its teledata
\emph{plus} telegate operations, so the baseline keeps its full action set and
is not handicapped by our scope; and Section~IV-E quantifies the regime in
which amortised gate teleportation genuinely wins.

\where{Section~I; Section~III opening; Section~IV-E; Section~VI; Table~III
caption.}

\begin{comment}
Comment 1.3. In Section~II-C the architecture is described as an abstract
black-box graph, yet the experiments rely on a specific grid-of-grids
structure. Moreover, the performance differences across circuit families
(e.g.\ degradation on star-structured GHZ) suggest sensitivity to the
underlying topology.
\end{comment}

\response
Two changes, one for each half of the comment.

First, the architectural assumptions are now stated exactly rather than left
as a black box: the router requires a connected intra-core subgraph per core
and a connected core graph, and queries the architecture only through three
distance tables. Degree, regularity, diameter, and the number and placement of
ports are \emph{not} assumed (Sections~I and~II-C).

Second, the evaluation now exercises that generality. Besides the two
grid-of-grids devices, all nine 64-qubit circuits are routed on four 27-qubit
IBM heavy-hex cores wired as a ring and as a star---cores with degrees
$\{1,2,3\}$, six degree-1 leaves, and diameter 12 against the $4\times4$
grid's 6, and, for the star, a hub that every inter-core route must cross.
Nothing is specialised for these devices: the layout reservation reads vertex
degree and remoteness, so it selects leaf qubits where it selected grid
corners. The reductions are $52.0\%$ (ring) and $41.5\%$ (star), with
\TeleSABRE{} again failing on Random where \dSABRE{} completes
(Section~IV-C).

On GHZ specifically: the regression you identified is gone (64-qubit GHZ is now
$-62.5\%$), and paragraph~(c) of the note above gives the full accounting.
The dominant cause on our side was not the score but the initial layout: the
submitted rule reserved escape capacity only in core~0, so GHZ landed with
four of six cores at $100\%$ occupancy. Reservation is now per core, and the
legality floor keeps ordinary routing from filling a core thereafter.

\where{Section~I; Section~II-C; Section~IV-C (``Topology''); Table~III.}

\begin{comment}
Comment 1.4. Section~III introduces a rich set of scoring functions and
mechanisms, but the presentation is complex and it is difficult to discern a
clear central principle; the formulation does not clarify why this particular
decomposition of terms is necessary.
\end{comment}

\response
Section~III is rewritten around Eq.~(3), which states the conceptual objective
in three labelled groups before any concrete term appears; the intra-core and
inter-core scores are then presented as two instances of it. All notation is
consolidated in Table~I with its defaults. Each term's role is now stated
where the term is introduced: why $c_{\mathrm{cap}}$ prices the immediate
landing core rather than the core holding the other operand; why the teleport
score, unlike Eq.~(6), carries no set-size normalisation (a teleport moves one
operand of one gate, and front-layer gates are qubit-disjoint, so exactly one
front-layer term ever changes); and why Eq.~(11) is an approximation, with its
two omissions stated explicitly. A running example (Figure~4 with Table~II)
isolates one factor per candidate and shows the three groups of Eq.~(3)
deciding among them. Finally, there are simply fewer mechanisms to follow than
before---three are removed---and Table~IV quantifies what each survivor
contributes.

\where{All of Section~III; Table~I; Figure~4; Table~II.}

\begin{comment}
Comment 1.5. Some comparisons are complicated by differences in communication
models and baseline capabilities. A more explicit discussion of the
equivalence (or lack thereof) between different EPR usage models would
strengthen the validity of the reported improvements.
\end{comment}

\response
The comparisons are now organised by how controlled they are, and labelled as
such.

\TeleSABRE{} is the primary head-to-head because it shares the teledata
setting, the architectures, and the local-SWAP accounting (Section~IV-A). The
\texttt{pytket-dqc} comparison is explicitly framed as a cost-model study
(Section~IV-E) and audited on the two assumptions that make its e-bit counts
incomparable with ours. A capacity audit shows that $4/6$, $2/6$, and at least
$6/9$ of its distributions provably cannot be materialised within the
device's communication-port capacity---25-qubit AE, for instance, requires 20
simultaneous link qubits at a core where the device provides two. Re-scoring
the same distributions under a bounded entanglement lifetime (at most three
gates per pair) turns the one suite it wins into a $59.6\%$ \dSABRE{}
advantage, which locates precisely where amortised gate teleportation is and
is not preferable. A DMapS comparison is added with its occupancy-preserving
two-EPR remote-SWAP model spelled out---because that primitive leaves each
chip's occupancy unchanged, DMapS never needs a free slot and maintains no
capacity condition during routing, which is a modelling choice rather than an
error but must be stated before the counts are read against ours. A new
``Device model'' paragraph in Section~V contrasts what each tool assumes of
the hardware, and notes that \dSABRE{}'s assumptions are the most conservative
of the four.

\where{Section~IV-A; Section~IV-E; Table~VI; Section~V.}

\section*{Reviewer 2}

\begin{comment}
Comment 2.1. The level of novelty appears moderate. The framework remains
firmly rooted in the \SABRE{}-style paradigm, and many contributions can be
viewed as heuristic refinements of existing routing methodologies rather than
fundamentally new algorithmic concepts.
\end{comment}

\response
This is the comment I answer by argument rather than by new experiments, so let
me be explicit about what is claimed as new.

\begin{enumerate}[leftmargin=*,itemsep=2pt]
\item \emph{A reformulation, not a re-weighting.} Prior \SABRE{}-style
distributed routers treat a teleport as an expensive edge. The claim here is
that a distributed \SABRE{} objective must additionally price the two state
changes a \emph{directional} teleport causes---port availability and
destination occupancy---and must price them inside one macro-action, because
staging, eviction, and the teleport itself succeed or fail together.
\item \emph{A completion guarantee.} Theorem~1 (termination with success under
checkable pre-routing conditions) is a property the compared \SABRE{}-style
distributed routers do not provide: \TeleSABRE{} can abort at its own deadlock
criterion, and DMapS's model has no capacity condition to maintain. It is also
what the completion result rests on ($45/45$ instances against $40/45$), and
it is validated at runtime rather than merely asserted---across 1074 recovery
transactions no asserted precondition failed.
\item \emph{Subtraction rather than addition.} Three mechanisms of the
original submission are removed, and the Conclusion claims the advance to be
the separation of ranking, legality, and recovery, not a count of terms.
\end{enumerate}

The \SABRE{} lineage itself is a deliberate design constraint, now argued
rather than assumed (Sections~I and~II-D): building on the decision rule that
production toolchains already run is what keeps compile time practical and the
output executable, and the research question is what that rule must become
when one chip becomes several. I would rather defend that position than claim
a novelty of paradigm the work does not have.

\where{Section~I (contributions); Section~II-D; Section~III-F; Section~V;
Section~VII.}

\begin{comment}
Comment 2.2. The paper relies predominantly on empirical evidence. There is
limited theoretical insight into why the proposed scoring formulation should
consistently outperform alternative approaches.
\end{comment}

\response
The revision adds the theory the problem admits, and is explicit about what
remains empirical.

New formal content: Lemmas~1--2 and Theorem~1 with proof (full version in the
supplement); pre-routing conditions that are verified before routing begins
(Algorithm~1, line~2), with the workload-side condition (Eq.~12) shown to be
loose---the device may be filled to $86\%$ for the 16-site cores used here and
$91$--$92\%$ for the 25-site ones, and the tightest instance in the paper has
32 free slots where 13 suffice; and a complexity analysis
(Section~III-H, Eq.~15, derived term by term in the supplement). The theory deliberately targets feasibility and termination,
which is the property whose absence shows up as baseline non-convergence, and
it is validated at runtime: none of the $1.08$ million scored candidates was
rejected at execution time, exactly as Lemma~2 predicts.

Score \emph{quality} is a different matter, and I do not claim a theorem for
it: which feasible progress-making move is cheapest is a heuristic-design
question. The revision supports it empirically on three axes---component
ablations (Table~IV), one-at-a-time parameter sweeps in which every default is
the best value swept on the 64-qubit suite and within $2.5\%$ of the best at
25 qubits, and a teleport-to-SWAP cost-ratio sweep from 10 to 100---and
Eq.~(11) is now labelled an approximation with its two omissions named rather
than presented as exact.

\where{Section~III-F; Section~III-H; Section~IV-D; supplementary material.}

\begin{comment}
Comment 2.3. Since the considered approaches employ different communication
primitives, direct comparisons based solely on EPR counts do not always
provide a fully controlled evaluation.
\end{comment}

\response
Please see the response to Comment~1.5 for the reorganised comparison
methodology, which is the substance of the answer. Three additions are
specific to control:

\begin{enumerate}[leftmargin=*,itemsep=2pt]
\item A \emph{shared-layout control} separates routing from allocation:
started from \TeleSABRE{}'s own initial layout, \dSABRE{}'s 64-qubit
geometric mean rises by $30.9\%$---and still improves on \TeleSABRE{} by
$33.4\%$ (Table~IV, Section~IV-D). The gain is therefore not an artefact of
our layout stage.
\item Reporting protocols are stated per tool's own convention (best of three
seeds for both routers, best of five for \texttt{pytket-dqc}), and, following
your concern about how much of the result rests on that convention,
Section~IV-B now also reports the median seed: the reductions become
$50.3/59.9/42.9\%$ against the best-of-three $50.8/62.0/49.1\%$. Full per-seed
counts are in the supplement.
\item Non-convergence is the tool's own verdict within a stated budget, is
reported explicitly rather than absorbed into a mean, and is excluded from
baseline-only geometric means, with the \dSABRE{}-only mean shown separately
(Section~IV-A, Table~III). We also decline to attribute the baseline's
non-convergence to capacity handling: it could equally reflect search
behaviour or implementation limits, and Section~IV-C says so.
\end{enumerate}

\where{Section~IV-A; Section~IV-B; Section~IV-D; Section~IV-E.}

\begin{comment}
Comment 2.4. The ablation indicates that the capacity-aware penalty
contributes a significant fraction of the gains, whereas some remaining
mechanisms provide comparatively modest improvements; a clearer discussion of
relative importance would help. In addition, incorporating latency, scheduling
effects, fidelity considerations, or more application-oriented workloads would
strengthen the practical relevance.
\end{comment}

\response
\textbf{(a) Relative importance.} The ablation is redesigned to answer this
directly. Capacity enters the design \emph{twice}, and Table~IV now separates
the two roles: removing the graded capacity score costs $+9.4\%$ EPR;
removing the legality floor alone costs $+0.2\%$; removing both costs
$+87.8\%$ \emph{and aborts two circuits}. The two overlap---the score's penalty band
contains the region the floor forbids, so either alone still keeps the router
off nearly-full cores on this suite---but only the floor supplies the
hypothesis Theorem~1 needs, a weight being outvotable where an admission test
is not. Section~IV-D now says exactly this rather than reporting a single
capacity number. Replacing the
BFS depth-ordered extended set with a topological one costs $+12.4\%$, and
tightening the floor to $\phi_1=3$ costs $+16.3\%$.

\textbf{(b) The modest mechanisms are gone.} Rather than defend components
whose contribution was small, I re-measured them and removed them: the
hop-gain reward (worth at most $+1.9\%$ in the original ablation, and net zero
or negative when re-measured across suites once congestion relief was
removed), the congestion-relief pass, and node decay. Every remaining
component now has a measured contribution in Table~IV.

\textbf{(c) Recovery is quantified as a safeguard, not an optimisation.} No
reported route at 25--64 qubits invokes it at all, while at 100, 200, and 360
qubits it retires 6, 10, and 66 gates (Table~V). This is stated as a bound on
its relevance, not as a contribution to the EPR gains.

\textbf{(d) Latency and fidelity.} Here I have to be straightforward: these
are outside the model this paper optimises, and I did not add a scheduling or
noise model rather than add a weak one. The scope is stated in Section~I; the
Conclusion now cautions explicitly that fewer EPR pairs does not by itself
imply a shorter makespan or a higher-fidelity result; and Section~VI sets out
what a scheduling and fidelity pass over the emitted traces would have to
account for---link and port conflicts, entanglement-generation rate, storage
lifetime, and hardware error rates. The routed traces are released with the
artefact precisely so that this can be built on them, and I would rather make
that a properly scoped follow-up than report a proxy metric whose modelling
assumptions would need as much defence as the result.

\textbf{(e) Workloads.} The 64-qubit suite adds QPEexact, QAOA, and a
$13\,040$-CX 16-bit Multiplier; the scalability study adds QPEexact alongside
QFT to 360 qubits; and the two heavy-hex networks vary the architecture rather
than the circuit. The evaluation now covers 45 circuit--architecture
instances.

\where{Tables~IV--V; Sections~IV-C and~IV-D; Section~VI; Section~VII.}

\section*{Reviewer 3}

\begin{comment}
Technical 3.1. The Abstract should be rewritten to state the motivation and
summarise the main contributions; in its current form it does not convey the
core idea.
\end{comment}

\response
The Abstract is rewritten around four things and nothing else: the motivation
(EPR pairs are the expensive resource, and finite ports and finite core
capacity constrain routing), the core idea in one sentence (each inter-core
move is scored as one macro-action---the EPR pair it spends, the port staging
it requires, the destination capacity it consumes, and the progress it buys),
the guarantee (\dSABRE{} is provably complete: every gate routed, no abort at
saturation, on any instance that passes a pre-routing check), and the
quantified results ($49$--$62\%$ geometric-mean EPR
reduction against \TeleSABRE{}, $45/45$ completions against $40/45$, and a
360-qubit QFT in $48$\,s). All acronyms are expanded. It is also deliberately
kept short, so that the core idea is not buried in a list of results.

\begin{comment}
Technical 3.2. The organization needs improvement: the Introduction states
that distributed compilation decomposes into three problem classes and begins
introducing them sequentially, but the discussion becomes overly detailed
before the third class is presented.
\end{comment}

\response
Fixed structurally. The three tasks---(i) allocation, (ii) choice of
communication primitive, (iii) routing and scheduling---are now presented
together in one compact enumerated passage, \emph{before} any of them is
discussed in depth. The following paragraphs are then explicitly scoped
(``A practical routing objective'', ``Scope of communication'', ``Scope of
architecture and comparison'', ``Contributions''), so the reader can see at
each point which of the three is being narrowed and which are being set aside.

\begin{comment}
Technical 3.3. In the Contributions, $\Delta_F$ and $\Delta_E$ are used
without definition, and ``comm port'' should be defined or expanded.
\end{comment}

\response
The contribution bullets are rewritten in words, with no undefined symbols:
they now speak of ``front-layer distance reduction plus DAG-depth-discounted
lookahead'' rather than $\Delta_F$ and $\widetilde{\Delta}_E$. All notation is
defined systematically in one place, Table~I, together with the default value
of every parameter. ``Communication port'' is spelled out and defined at first
use in Section~II-C: an ordinary physical slot that also terminates an
inter-core link, and that may hold a data qubit when it is not being used for
communication---which is exactly why it may have to be evacuated before a
teleport, and hence why the score carries an availability term at all.

\begin{comment}
Technical 3.4. In the Background, ``Together, CX and arbitrary 1Q gates form a
universal gate set'' should specify the 1Q gate sets.
\end{comment}

\response
The statement is made precise (Section~II-A): CX together with \emph{arbitrary}
single-qubit unitaries is \emph{exactly} universal, in the sense that it
decomposes any $n$-qubit unitary (Barenco et al., now cited at that point),
whereas a finite set such as $\{\mathrm{CX},H,T\}$ is only
\emph{approximately} universal. Since hardware provides a finite set, the
paragraph also states what the benchmarks actually use: all circuits are
transpiled into IBM's native basis $\{\mathrm{CX},R_z,\sqrt{X},X\}$ at Qiskit
optimisation level 3.

\begin{comment}
Technical 3.5. The claim that \dSABRE{} is slower than \TeleSABRE{} for
circuit compilation requires further justification.
\end{comment}

\response
Rather than justify that claim, I have removed its cause: it no longer holds
for the revised implementation. Paragraph~(d) of the note above gives the
engineering account---distance queries are composed through Eq.~(1) and
memoised instead of being answered by per-iteration graph searches, the DAG
bookkeeping is incremental, and the removed mechanisms cut per-iteration
work---and confirms that the rewrite is output-identical to the previous
implementation over 207 routing passes.

The measured picture, end-to-end and including layout generation and all nine
routing passes, is a geometric mean of $1.9$\,s for \dSABRE{} against $4.7$\,s
for \TeleSABRE{} on the 64-qubit suite, and $23.7$\,s against $136.3$\,s on
the $13\,040$-CX Multiplier. Section~IV-C reports this without
over-claiming: the gap is not uniform (on QNN the two are within $0.2$\,s),
the two tools run different search procedures, so these are end-to-end
practical costs rather than an operation-level comparison, and \dSABRE{}
remains an interpreted implementation competing with a compiled one---the
machine, the languages, and the tool versions are now stated in Section~IV-A.
Within one architecture, compile time is near-linear in gate count (log--log
exponent $0.97$, $R^2=0.95$ over the nine 64-qubit circuits). Section~IV-C
states that each such fit holds the device fixed, and the supplement fits the
scalability series separately---exponent $2.17$, because there the device grows
with the circuit, which Eq.~(15) charges through its $K^2$ term---so the
near-linear reading is not mistaken for a claim about scaling the machine.

\begin{comment}
Technical 3.6. Fig.~1 nicely explains the DAG representation and Fig.~3
presents the workflow clearly, but the caption of Fig.~4 is unclear and should
be revised for better readability.
\end{comment}

\response
The Figure~4 caption is rewritten as a guided walk-through rather than a list
of graphical conventions: it names the pending cross-core gate and where its
two operands sit, states that the router scores the three links leaving the
source core and takes the cheapest by Eq.~(7), and then walks the two steps of
the macro-action---intra-core SWAPs to the staging slot (the count on the
solid arrow), then the teleport across the link (the dashed arc, annotated
with the resulting score)---before explaining each remaining annotation,
including which core is nearly full and which port is already occupied. The
candidates are labelled A, B, C to match Table~II row by row. One symbol in
the figure disagreed with the text ($f_{\mathrm{nxt}}$ against
$f_{c_{\mathrm{next}}}$) and has been unified. The Figure~3 caption is
likewise clarified: it now names all three conditions that invoke the recovery
transaction and states which of the three is drawn.

\begin{comment}
Editorial 3.1. The acronym DSABRE is written inconsistently throughout (both
lowercase and uppercase initial letters in title, abstract, and introduction).
\end{comment}

\response
Corrected. The name is set by a single macro and now reads ``\dSABRE{}''
uniformly---leading lowercase ``d'', small-capital stem---in the title,
abstract, body, figures, table headers, and captions.

\begin{comment}
Editorial 3.2. Abbreviations should be expanded at first appearance (EPR, DAG,
MQT Bench and similar), particularly in the Abstract.
\end{comment}

\response
Done. Einstein--Podolsky--Rosen (EPR), directed acyclic graph (DAG), and
Munich Quantum Toolkit (MQT) Bench are expanded at first occurrence in the
Abstract, and DAG is re-introduced with its definition at first use in the
body (Section~II-A), where the front layer is defined with it.

\begin{comment}
Overall assessment. Several notations and abbreviations are introduced without
definition. I strongly recommend defining all notation systematically in one
place. A more structured presentation and substantial reorganisation would
greatly improve readability.
\end{comment}

\response
Both recommendations are followed structurally rather than cosmetically.
Table~I is now the single notation reference for the whole paper, listing every
symbol with its meaning and, for each parameter, the default used in every
experiment. The reorganisation is described in the Summary of Major Changes
above; the parts that most affect readability are the rewritten Section~III
built around one objective, the compact three-task passage in the
Introduction, the scoped introduction paragraphs, and the running example that
shows the score deciding an actual case before the reader meets the general
algorithm. Section~IV-A also now states the machine, implementation language,
and tool versions, and Section~IV-E names the exact \texttt{pytket-dqc}
distributor used, so that the experiments can be reproduced from the text.

\section*{Area Editor (Prof.\ O'Connor)}

Thank you for the clear synthesis. Mapping the four requirements to the
revision:

\begin{enumerate}[leftmargin=*,itemsep=2pt]
\item \textbf{Core novelty and unifying design principle} --- the macro-action
formulation of Eq.~(3), the separation of ranking from legality from recovery,
and Theorem~1 (responses 1.1 and 2.1); three mechanisms removed rather than
added to.
\item \textbf{Justification of key design choices and how they work together}
--- the scope paragraphs of Section~I, the measured teledata-versus-telegate
analysis, the running example, and the redesigned ablation that separates the
capacity score's efficiency role from the legality condition the completion
proof rests on
(responses 1.2, 1.4, 2.4).
\item \textbf{Exposition and notation} --- rewritten abstract and
introduction, Table~I as the single notation reference, rewritten figure
captions, consistent naming, and the reproducibility details now in
Section~IV-A (responses 3.1--3.6 and the editorial items).
\item \textbf{Experimental comparisons across communication primitives and EPR
cost models} --- \TeleSABRE{} as the controlled like-for-like baseline, the
\texttt{pytket-dqc} capacity and entanglement-lifetime audits, the DMapS
comparison with its device model stated, and the ``Device model'' paragraph of
Section~V (responses 1.5 and 2.3).
\end{enumerate}

I would also draw your attention to paragraph~(b) of the note above: while
regenerating the results I found that the original submission had passed a
misspelled parameter to the \TeleSABRE{} baseline. It is corrected, the
baseline was re-run unchanged, and the correction makes the baseline stronger
on the circuits that dominate the geometric means.

\section*{Associate Editor (Prof.\ Chattopadhyay)}

Thank you. The three concerns are addressed respectively in response~2.1
(novelty, together with the explicit statement of what is and is not claimed
as new); in the broadened evaluation (nine 64-qubit circuit families, four
architectures, two circuit families scaled to 360 logical qubits, 45
circuit--architecture instances, and a redesigned ablation); and in the
reorganised exposition (the Reviewer~3 responses, Table~I, and the rewritten
Sections~I and~III).

\bigskip\noindent
I believe the revision addresses all of the comments, and I hope the
manuscript is now suitable for publication in IEEE TCAD. I would be glad to
make further changes.

\bigskip\noindent
Sincerely,\\
Sanjiang Li

\end{document}
